\documentclass{situationist-preprint}

\title{The Incredulous Stare as an Argument}
\msid{TSM-08-b}
\msfamily{TSM-08}
\msdiscipline{Time, Space and Modality}
\mstype{Companion manuscript}
\mskeywords{modal realism; incredulous stare; ontological cost; common sense; Lewis}

\begin{document}
\maketitle

\begin{abstract}
Lewis conceded that the commonest response to modal realism is an incredulous stare, and replied that a stare is not an argument. We argue that it can be made into one. The stare registers a cost: modal realism contradicts a firm belief that there is nothing concrete beyond the one world we inhabit. A cost of this kind can be outweighed only by theoretical benefit, and TSM-08-a shows that concrete worlds deliver no benefit that the ersatz account lacks. Reconstructed in this way, the stare is a valid argument whose premises Lewis himself accepted, and it supports the conclusion that other possible worlds do not exist concretely.
\end{abstract}

\section{From stare to argument}

\citet[\S2.8]{Lewis1986} grants that his view disagrees with common opinion, that this is a serious cost, and that the cost is not met by refutation, since the incredulous stare contains no premise he could deny. He proposes to pay it out of the theoretical benefits of the plurality of worlds. The stare is often dismissed as a refusal to engage \citep{Lycan1991}. We take Lewis's own accounting seriously instead, and state it as follows.

\begin{principle}[Conservatism]\label{p:cons}
A theory that contradicts a firmly held belief incurs a cost, which is acceptable only if the theory does work that no theory consistent with the belief does at lower cost.
\end{principle}

\begin{assumption}\label{a:belief}
It is firmly held that there are no concrete donkeys, planets or persons that stand in no spatiotemporal relation to us.
\end{assumption}

\begin{proposition}\label{p:stare}
Given Principle~\ref{p:cons}, Assumption~\ref{a:belief} and Theorem~4.2 of TSM-08-a, belief in concrete worlds other than the actual one is not warranted.
\end{proposition}

\begin{proof}
Modal realism contradicts the belief of Assumption~\ref{a:belief}, since it holds that there are concrete talking donkeys at other worlds, isolated from ours. By Principle~\ref{p:cons} the cost is acceptable only if modal realism does work that no theory consistent with the belief does. The ersatz account of TSM-07 is consistent with it, and by Theorem~4.2 of TSM-08-a it does all the work that concrete worlds do. The cost is therefore unpaid.
\end{proof}

The proposition does not make the stare redundant. Principle~P1 of TSM-08-a speaks of ontological cost in general; the stare identifies the particular cost incurred and shows that it is high. A theory that posits isolated concrete donkeys contradicts not a philosophical prejudice but the content of ordinary beliefs about what there is. Principle~\ref{p:cons} and Principle~P1 of TSM-08-a converge: the first counts the conflict with opinion, the second the ontology, and both demand the same payment.

\section{Objections}

\begin{objection}
Common opinion about what exists is restricted to what is actual, so modal realism does not contradict it: the ordinary claim is that there are no talking donkeys \emph{actually}, which Lewis accepts.
\end{objection}

\begin{reply}
The reply requires that the ordinary quantifier be restricted to our world. But the ordinary believer who denies that there are talking donkeys anywhere does not intend the restriction, and would deny the unrestricted claim when it is put to her \citep[compare][]{Stalnaker1976}. Lewis's reading of the belief is a reinterpretation, and reinterpretation is itself a cost.
\end{reply}

\begin{objection}
Common sense has been overturned before, by physics. Its verdicts on recondite matters carry little weight.
\end{objection}

\begin{reply}
Physics overturned common opinion by evidence and by explanatory success. Concrete worlds are, by hypothesis, inaccessible to evidence, and TSM-08-a shows that their explanatory success is not distinctive. The precedent therefore supports the stare: common opinion yields where something is gained by yielding, and here nothing is.
\end{reply}

\section{Conclusion}

The incredulous stare is the informal expression of a valid argument. Its premises are a principle of conservatism that Lewis accepted, a belief that he acknowledged is widely held, and the result of TSM-08-a that concrete worlds do no distinctive work. Other possible worlds do not exist concretely.

\begin{thebibliography}{99}\small
\bibitem[Lewis(1986)]{Lewis1986} Lewis, D. (1986). \emph{On the Plurality of Worlds}. Blackwell.
\bibitem[Lycan(1991)]{Lycan1991} Lycan, W.~G. (1991). Two---no, three---concepts of possible worlds. \emph{Proceedings of the Aristotelian Society}, 91, 215--227.
\bibitem[Stalnaker(1976)]{Stalnaker1976} Stalnaker, R.~C. (1976). Possible worlds. \emph{No\^us}, 10(1), 65--75.
\end{thebibliography}

\end{document}
